\cgeentry{cge_parameter}{User parameters: values, assignments and data}

\cgesec{Description}

\cmd{cge_parameter} adds a parameter whose values the user gives, with no
index, one or two, over model sets or user sets (\manref{cge_set}). User
variables (\manref{cge_variable}) and equations (\manref{cge_equation}) refer
to it by name. The values given are its base values; \manref{cge_shock}
changes them in a scenario, \manref{cge_optimize} may take one of its elements
as an instrument, continuation and dynamic runs carry it like the model's own
parameters, and \manref{cge_gams} writes it into the GAMS program. From
version 0.13.0 its value may be an expression, evaluated element by element
over the elements that meet a condition, as GAMS's assignments, so that
parameters are calibrated in the language, and data in wide form, as a SAM
in memory, give a parameter with two indexes.

Do not confuse \cmd{cge_parameter}, which defines a user parameter, with
\manref{cge_parameters}, which lists every parameter of the model, user
parameters included, as a dataset.

\cgesec{Quick start}

A scalar target, 1.1
\begin{stcode}
cge_parameter hgoal = 1.1
\end{stcode}
Weights by commodity, the elements not given 0
\begin{stcode}
cge_parameter w(c), values(com_agr:0.4 com_man:0.3 com_ser:0.2 com_hlt:0.1)
\end{stcode}
Shares by commodity and household, from the data in memory
\begin{stcode}
cge_parameter sh(c,h), from(share) labels(commodity household)
\end{stcode}
A scenario with a higher target
\begin{stcode}
cge_shock hgoal, value(1.2)
\end{stcode}
Assignments: shares of intermediate input by activity, and half the tariff
rates where there are tariffs
\begin{stcode}
cge_parameter qint0(c,a) = base(QINT(c,a))
cge_parameter share(c,a) = qint0(c,a) / sum(c2, qint0(c2,a))
cge_parameter tmr(c)$(tm(c) > 0) = tm(c) / 2
\end{stcode}
A SAM in memory, one variable for each account, as a parameter
\begin{stcode}
cge_parameter sam(ac,ac2), from(act_agr-row) labels(account)
\end{stcode}

\cgesec{Syntax}

A parameter assigned an expression (a number, or from version 0.13.0 any
expression of parameters)
\begin{stsyntax}
\cmdb{cge_parameter} \argu{name}[\cmd{(}\argu{i}[\cmd{,}\argu{j}]\cmd{)}][\cmd{$(}\argu{condition}\cmd{)}] \cmd{=} \argu{expression}
[\texttt{,} \optarg{def}{ault}{\#}]
\end{stsyntax}
Values by element
\begin{stsyntax}
\cmdb{cge_parameter} \argu{name}\cmd{(}\argu{i}[\cmd{,}\argu{j}]\cmd{)}\texttt{,}
\optarg{val}{ues}{spec} [\optarg{def}{ault}{\#}]
\end{stsyntax}
Values from the data in memory
\begin{stsyntax}
\cmdb{cge_parameter} \argu{name}\cmd{(}\argu{i}[\cmd{,}\argu{j}]\cmd{)}\texttt{,}
\optarg{}{from}{varname} \optarg{lab}{els}{varlist} [\optarg{def}{ault}{\#}]
\end{stsyntax}
Values from data in wide form (from version 0.13.0)
\begin{stsyntax}
\cmdb{cge_parameter} \argu{name}\cmd{(}\argu{i}\cmd{,}\argu{j}\cmd{)}\texttt{,}
\optarg{}{from}{varlist} \optarg{lab}{els}{varname} [\optarg{def}{ault}{\#}]
\end{stsyntax}
Display or remove the user parameters
\begin{stsyntax}
\cmdb{cge_parameter} \texttt{,} \opt{}{list}\\
\cmdb{cge_parameter} \texttt{,} \opt{}{clear}
\end{stsyntax}
\argu{i} and \argu{j} are indexes of model sets (\cmd{a c f h d i n}) or of
user sets; in a standalone model (\manref{cge_model}), of user sets only.
\argu{name} follows the rules of \manref{cge_set}. \argu{expression} is
written in the language of \manref{cge_equation}: numbers, user parameters,
the model's parameters, base values of variables (\cmd{base(}\argu{X}\cmd{)}),
sums, conditions and functions, with the indexes \argu{i} and \argu{j}.

\begin{opttable}
\optgroup{Main}
\optarg{val}{ues}{spec} & values by element, \argu{element}\cmd{:}\argu{\#} or \argu{e1}\cmd{.}\argu{e2}\cmd{:}\argu{\#}\\
\optarg{}{from}{varname} & a variable of the data in memory with the values\\
\optarg{lab}{els}{varlist} & the string variables with the element labels, one for each index\\
\optarg{}{from}{varlist} & in wide form, one variable for each element of the second index\\
\optarg{def}{ault}{\#} & value of the elements not given; default \texttt{0}\\
\end{opttable}

\cgesec{Options}

\begin{optdesc}
\item[\optarg{val}{ues}{spec}] gives values by element, as
\argu{element}\cmd{:}\argu{\#} pairs separated by spaces, for example
\cmd{values(com_agr:0.3 com_man:0.5)}; for a parameter with two indexes,
\argu{e1}\cmd{.}\argu{e2}\cmd{:}\argu{\#}.

\item[\optarg{}{from}{varname}] takes the values from the numeric variable
\argu{varname} of the data in memory, one observation per element.
Observations with a missing value or an empty label are skipped.

\item[\optarg{lab}{els}{varlist}] names the string variables that hold the
labels of the elements, one for each index, with \opt{}{from()}.

\item[\optarg{}{from}{varlist} \optarg{lab}{els}{varname}] (from version
0.13.0) read data in wide form for a parameter with two indexes, as a SAM in
memory: the string variable \argu{varname} labels the elements of the first
index, one observation each, and each variable of \argu{varlist},
abbreviations and ranges allowed, holds the values for the element of the
second index whose label is its name.

\item[\optarg{def}{ault}{\#}] gives the value of the elements that
\opt{val}{ues()} or \opt{}{from()} leaves out; the default is 0.

\item[\opt{}{list}] displays the user parameters, with the range of their
current and base values.

\item[\opt{}{clear}] removes every user parameter. A parameter that a user
variable or equation uses cannot be removed until they are.
\end{optdesc}

\cgesec{Remarks and examples}

The SAM with a health sector is mapped and calibrated:
\stlog{par_setup}

\cgesubsec{Example 1: a target set by a parameter}

A target for government purchases of health services, 1.1 times the base,
which a user equation enforces with the scaling factor of government
consumption adjusting:
\stlog{par_scalar}
The base does not satisfy the equation, which therefore acts as a shock.

\cgesubsec{Example 2: a shock to a parameter}

\manref{cge_shock} changes a user parameter like a parameter of the model:
\stlog{par_shock}
The scaling factor of government consumption raises every government purchase
in proportion, so real government consumption is 10 and then 20 per cent above
the base; investment falls by about as much.

\cgesubsec{Example 3: values by element}

Weights by commodity, for a price index of the user's own:
\stlog{par_values}
From the data in memory, with the labels in a string variable; the elements
without an observation take \opt{def}{ault()}, 0:
\stlog{par_from}
With two indexes, here commodity and household:
\stlog{par_two}

\cgesubsec{Example 4: the parameters as data}

\manref{cge_parameters} lists user parameters after those of the model:
\stlog{par_data}

\cgesubsec{Example 5: a parameter as an instrument}

The health purchase target that leaves government savings at 1.5, against 1
in the base:
\stlog{par_opt}

\cgesubsec{Example 6: assignments}

From version 0.13.0 a parameter's value is an expression, evaluated for each
element. The shares of each commodity in the intermediate input of each
activity, from the base values of the variables; half the import tariff rate,
where there is a tariff, and $-1$ elsewhere, then 0 where it is $-1$, by a
second assignment:
\stlog{par_assign}
An assignment is evaluated when it is made: \cmd{share} does not change if
\cmd{QINT} does, and \cmd{tmr} does not follow a later shock to \cmd{tm}.

\cgesec{Stored results}

\cmd{cge_parameter} stores the following in \cmd{r()}:
\begin{storedres}
\resgroup{Scalars}
\res{r(N_elements)} & number of values of the parameter defined\\
\res{r(N_parameters)} & number of user parameters (\opt{}{list})\\
\resgroup{Macros}
\res{r(name)} & name of the parameter defined\\
\end{storedres}

\cgesec{Methods and formulas}

The values of all user parameters are held, row by row, in one vector of the
model's parameters, so every simulation stores them with its other
parameters, \cmd{cge_shock, reset} returns them to their base values, and
continuation in $n$ steps moves them as $\theta_k = \theta^0 + (k/n)(\theta^1
- \theta^0)$ like the others. In an expression, a reference to a parameter
element compiles to its position in that vector; its derivative is zero. When
the model is calibrated again, the parameters keep their base values. In the
GAMS program, \cmd{cge_user.inc} declares each parameter as
\cmd{UP_}\argu{name}, with its base values, and the scenario file changes it
by the ratio of its value to the base, as for the model's parameters.

An assignment writes its expression out for each element of its domain that
meets its condition at the base, compiles it as an equation's side is
compiled and evaluates it with the base values of the parameters and
variables; an expression that refers to a variable, rather than to its base
value, is an error, as is an element whose value is missing, such as a
division by zero.

\cgesec{Also see}

\begin{alsosee}
\seeentry{cge_set}{User sets of elements, for indexes and sums}
\seeentry{cge_variable}{User variables, defined by an expression or determined by user equations}
\seeentry{cge_equation}{User equations: new relations among the model's variables}
\seeentry{cge_shock}{Change a parameter or a fixed variable}
\seeentry{cge_parameters}{The calibrated parameters as a dataset}
\seechapter{chap:user}{User variables and equations}
\end{alsosee}
